\chapter{आनत समष्टियाँ}\label{ch:b2:affine}

सदिश समष्टियों में एक विशेषाधिकार प्राप्त बिंदु होता है --- मूल बिंदु ---
जो ज्यामिति को नहीं चाहिए। \emph{\omterm{def:b2:affine:def}{आनत समष्टि}} वह सदिश समष्टि है जो अपना मूल
बिंदु भूल चुकी है: बिंदु और सदिश अलग-अलग जातियाँ बन जाते हैं, जिन्हें
स्थानांतरण जोड़ता है। यह छोटा अध्याय वह शब्दकोश बनाता है (बिंदु, \omterm{def:b2:affine:barycenter}{बैरिकेंद्र},
\omterm{def:b2:affine:subspace}{आनत} उपसमष्टियाँ और प्रतिचित्रण), उत्तलता का \omterm{def:b2:affine:subspace}{आनत} दृष्टिकोण, और आगे के
ज्यामिति अध्यायों में काम आने वाले वर्गीकरण-औज़ार।

\section{बिंदु और सदिश}

\begin{definition}\label{def:b2:affine:def}
किसी वास्तविक सदिश समष्टि $E$ से दिशित \emph{आनत समष्टि}\index{आनत समष्टि}
कोई अरिक्त समुच्चय $\mathcal{E}$ है जिस पर ऐसा प्रतिचित्रण $(A, B) \mapsto \vect{AB} \in E$ हो कि
\[
\vect{AB} + \vect{BC} = \vect{AC}
\quad \text{(शाल)},
\qquad
\text{प्रत्येक } A,\ B \mapsto \vect{AB} \text{ के लिए एकैकी आच्छादक }
\mathcal{E} \to E .
\]
और $\vect{AB} = u$ वाले अद्वितीय बिंदु को $B = A + u$ लिखा जाता है। $\mathcal{E}$ की \emph{विमा} $\dim E$
है। हर सदिश समष्टि अपने ऊपर \omterm{def:b2:affine:subspace}{आनत} समष्टि है ($\vect{AB} = B - A$); और मूल बिंदु $O \in \mathcal{E}$ का हर
चुनाव $M \mapsto \vect{OM}$ के द्वारा $\mathcal{E}$ को $E$ से पहचान देता है।
\end{definition}

\begin{example}[बिना स्वाभाविक मूल बिंदु की एक आनत समष्टि]\label{ex:b2:affine:planeexample}
हल-समतल $\mathcal E = \{(x, y, z) \in \R^3 : x + y
+ z = 1\}$ सदिश उपसमष्टि नहीं है ($0 \notin \mathcal E$), परंतु वह $E = \{x + y + z =
0\}$ से दिशित \omterm{def:b2:affine:def}{आनत समष्टि}
है: $A, B \in \mathcal E$ के लिए अंतर $\vect{AB} = B
- A$ $E$ में उतरता है (योग कट जाते हैं), शाल $\R^3$ से
विरासत में मिलती है, और $B \mapsto \vect{AB}$ $E$ पर एकैकी आच्छादक है। $\mathcal E$ का कोई बिंदु
विशिष्ट नहीं है --- ``मूल बिंदु'' का कोई भी चुनाव $O \in \mathcal E$ उतना ही अच्छा है,
और सारी पहचानें $M \mapsto \vect{OM}$ स्थानांतरणों से भिन्न होती हैं। यही विशिष्ट स्थिति
है: असमांगी रैखिक समस्याओं (रैखिक निकाय, \cref{ch:b2:diffeq} में रैखिक अवकल समीकरण) के
हल-समुच्चय \omterm{def:b2:affine:subspace}{आनत} होते हैं, रैखिक कभी नहीं; और नारा ``विशिष्ट हल जमा केंद्रक''
ठीक अगली परिभाषा का कथन $\mathcal F = A + F$ है।
\end{example}

\begin{definition}[बैरिकेंद्र]\label{def:b2:affine:barycenter}
मान लीजिए $(A_i, \lambda_i)_{i \leq k}$ भारित बिंदु हैं जहाँ $\sum\lambda_i \neq 0$। \emph{बैरिकेंद्र}\index{बैरिकेंद्र}
$G
= \operatorname{bar}\bigl((A_i, \lambda_i)\bigr)$ वह अद्वितीय बिंदु है जिसके लिए
\[
\sum_i \lambda_i\, \vect{GA_i} = 0
\qquad\text{समतुल्य रूप से}\qquad
\vect{OG} = \frac{1}{\sum\lambda_i}\sum_i \lambda_i\,\vect{OA_i}
\quad (\text{किसी भी } O).
\]
बैरिकेंद्र \emph{साहचर्यी} होते हैं (बिंदुओं के उपसमूहों के स्थान पर उनका
आंशिक बैरिकेंद्र जोड़े गए भार के साथ रखा जा सकता है) और सारे भारों के मापन
बदलने पर अपरिवर्त्य रहते हैं।
\end{definition}

\begin{proof}[अस्तित्व और सूत्रों की उपपत्ति]
$O$ स्थिर कीजिए और $s = \sum_i\lambda_i \neq 0$ लिखिए। शाल से
\[
\sum_i\lambda_i\,\vect{GA_i} = 0
\iff \sum_i\lambda_i\bigl(\vect{GO} + \vect{OA_i}\bigr) = 0
\iff s\,\vect{OG} = \sum_i\lambda_i\,\vect{OA_i},
\]
जो $G = O + \frac1s\sum_i\lambda_i\vect{OA_i}$ को अद्वितीय रूप से निर्धारित कर देता है। \emph{$O$ से स्वतंत्रता:}
किसी अन्य मूल बिंदु $O'$ के लिए
\[
\frac1s\sum_i\lambda_i\,\vect{O'A_i}
= \frac1s\sum_i\lambda_i\bigl(\vect{O'O} + \vect{OA_i}\bigr)
= \vect{O'O} + \frac1s\sum_i\lambda_i\,\vect{OA_i}
= \vect{O'G} :
\]
अर्थात् वही बिंदु $G$। \emph{साहचर्यता:} सूचकांक समुच्चय को $s_I = \sum_{i\in I}\lambda_i \neq 0$ के साथ
$I \sqcup J$ में बाँटिए, और $G_I$ को $(A_i, \lambda_i)_{i\in
I}$ का \omterm{def:b2:affine:barycenter}{बैरिकेंद्र} लीजिए, ताकि $\sum_{i\in I}\lambda_i\vect{OA_i} =
s_I\,\vect{OG_I}$। तब
\[
s\,\vect{OG}
= \sum_{i\in I}\lambda_i\vect{OA_i} +
\sum_{j\in J}\lambda_j\vect{OA_j}
= s_I\,\vect{OG_I} + \sum_{j\in J}\lambda_j\vect{OA_j} :
\]
अर्थात् $G$ $(G_I, s_I)$ तथा $(A_j,
\lambda_j)_{j\in J}$ का \omterm{def:b2:affine:barycenter}{बैरिकेंद्र} है, जैसा दावा किया गया।
\emph{मापन-परिवर्तन:} हर $\lambda_i$ के स्थान पर $t\lambda_i$ ($t \neq 0$) रखने से $s$ और
भारित योग दोनों $t$ से गुणा हो जाते हैं, जिससे $\vect{OG}$ अपरिवर्तित रहता है।
\end{proof}

\begin{remark}[सामान्य भूलें]
परिभाषा के चारों ओर दो जाल हैं। पहला, यदि भारों का योग \emph{शून्य} हो तो
कोई \omterm{def:b2:affine:barycenter}{बैरिकेंद्र} नहीं होता: तब प्रतिचित्रण $O \mapsto
\sum\lambda_i\vect{OA_i}$ $O$ से स्वतंत्र हो जाता है
और कोई \emph{सदिश} परिभाषित करता है, बिंदु नहीं --- उदाहरण के लिए $(A, -1;\ B, 1)$ $\vect{AB}$
को संकेतित करता है। किसी परिकलन से दोनों में से क्या निकल रहा है, इसका
हिसाब रखना बैरिकेंद्रीय स्वच्छता का आधा हिस्सा है। दूसरा, भार केवल किसी
साझे अशून्य गुणक तक ही अर्थपूर्ण होते हैं; और ``$G$ के निर्देशांक $\lambda_1, \dots, \lambda_k$ हैं''
जैसे सूत्र किसी सामान्यीकरण (प्रायः $\sum\lambda_i = 1$) को मान लेते हैं, तथा सामान्यीकरण
भूल जाना किसी चित्र पर ग़लत अनुपातों का मानक स्रोत है।
\end{remark}

\begin{definition}[आनत उपसमष्टियाँ; आनत प्रतिचित्रण]\label{def:b2:affine:subspace}
\emph{आनत उपसमष्टि} ऐसा समुच्चय $\mathcal{F} = A + F = \{A + u :
u \in F\}$ है जिसके लिए $F$ कोई सदिश उपसमष्टि हो
(उसकी \emph{दिशा}); और समतुल्य रूप से, बैरिकेंद्रों के अंतर्गत स्थायी कोई
अरिक्त समुच्चय। $\R^n$ की आनत उपसमष्टियाँ ठीक रैखिक निकायों $MX = B$ के हल-समुच्चय
हैं (प्रथम वर्ष: विशिष्ट हल जमा केंद्रक)। कोई प्रतिचित्रण $f \colon
\mathcal{E} \to \mathcal{E}'$
\emph{आनत}\index{आनत प्रतिचित्रण} कहलाता है जब वह \omterm{def:b2:affine:barycenter}{बैरिकेंद्र} सुरक्षित रखे
--- समतुल्य रूप से, जब किसी (अद्वितीय) रैखिक प्रतिचित्रण $\varphi = \vec f$, अर्थात्
\emph{रैखिक भाग}, के लिए
\[
f(A + u) = f(A) + \varphi(u)
\]
हो। $\R^n$ के आनत प्रतिचित्रण: $X \mapsto MX + C$. संयोजन आनत होते हैं और उनके रैखिक भाग
संयोजित; और $f$ एकैकी आच्छादक है तभी जब $\vec f$ हो।
\end{definition}

\begin{proof}[प्रतिचित्रणों के लिए तुल्यता की उपपत्ति]
यदि $f(A + u) = f(A) + \varphi(u)$: तो $(A_i,
\lambda_i)$ के किसी \omterm{def:b2:affine:barycenter}{बैरिकेंद्र} $G$ के लिए हर बिंदु को $A$ से खोलने पर
$f(G) = f(A) +
\varphi(\vect{AG})$ और $\varphi(\vect{AG}) =
\frac{\sum\lambda_i\varphi(\vect{AA_i})}{\sum\lambda_i}$: अर्थात् $f(G)$ प्रतिबिंबों का \omterm{def:b2:affine:barycenter}{बैरिकेंद्र} है। विलोमतः $A$ स्थिर
कीजिए और $\varphi(u) = \vect{f(A)\,f(A + u)}$ परिभाषित कीजिए। \emph{समघातता:} प्रत्येक वास्तविक $t$ के लिए
$A + tu =
\operatorname{bar}\bigl(A, 1-t;\ A + u, t\bigr)$, अतः (परिकल्पना के अनुसार स्वेच्छ वास्तविक भारों वाले) बैरिकेंद्रों का
सुरक्षित रहना सीधे $\varphi(tu) = t\,\varphi(u)$ दे देता है। \emph{योज्यता:} $A + u + v = \operatorname{bar}\bigl(A + 2u,
\tfrac12;\ A + 2v, \tfrac12\bigr)$, अतः समघातता का
उपयोग करके $\varphi(u + v) =
\tfrac12\varphi(2u) + \tfrac12\varphi(2v) = \varphi(u) +
\varphi(v)$। इस प्रकार $\varphi$ रैखिक है।
\end{proof}

\begin{remark}
उपपत्ति ने \emph{स्वेच्छ वास्तविक} भारों वाले \omterm{def:b2:affine:barycenter}{बैरिकेंद्र} काम में लिए:
समघातता वाला पग $t$ को $\intcc01$ के बाहर ले जाता है। यदि किसी प्रतिचित्रण के बारे
में केवल यह माना जाए कि वह अऋणात्मक भारों वाले \omterm{def:b2:affine:barycenter}{बैरिकेंद्र} सुरक्षित रखता है
--- अर्थात् मध्यबिंदु और खंड --- तो सदिश प्रतिचित्रण की रैखिकता मुफ़्त में
नहीं मिलती: केवल $\Q$-रैखिकता मिलती है, और निष्कर्ष के लिए कोई संततता
परिकल्पना चाहिए, ठीक जैसा \cref{exo:b2:affine:5} में। ``सारे \omterm{def:b2:affine:barycenter}{बैरिकेंद्र} सुरक्षित रखता है''
और ``उत्तल संचय सुरक्षित रखता है'' में भेद करना \omterm{def:b2:affine:subspace}{आनत} शब्दावली की एक छोटी पर
सच्ची सूक्ष्मता है।
\end{remark}

\begin{example}[क्लासिक बैरिकेंद्रीय ज्यामिति]\label{ex:b2:affine:median}
त्रिभुज $ABC$ का गुरुत्व केंद्र \omterm{def:b2:affine:barycenter}{बैरिकेंद्र} $G =
\operatorname{bar}(A,1; B,1; C,1)$ है। मध्यबिंदु $A' = \operatorname{bar}(B, 1; C, 1)$ के साथ
साहचर्यता दिखाती है
\[
G = \operatorname{bar}(A, 1;\ A', 2) :
\]
अर्थात् $G$ माध्यिका $AA'$ पर उसके दो-तिहाई भाग पर पड़ता है --- और शेष दोनों
माध्यिकाओं के लिए भी वैसा ही: इस प्रकार तीनों माध्यिकाएँ संगामी हैं, और वह
भी बैरिकेंद्रीय कलन की एक ही पंक्ति में।
\end{example}

\begin{example}[किसी चतुर्भुज की द्वि-माध्यिकाएँ]\label{ex:b2:affine:bimedians}
मान लीजिए $ABCD$ कोई भी चतुर्भुज है (समतलीय हो या न हो!) और उसकी
\emph{द्वि-माध्यिकाएँ} लीजिए: अर्थात् आमने-सामने की भुजाओं के मध्यबिंदुओं
को जोड़ने वाले खंड $M_{AB}M_{CD}$ और $M_{BC}M_{DA}$। $(A,1; B,1; C,1; D,1)$ का \omterm{def:b2:affine:barycenter}{बैरिकेंद्र} $G$ लीजिए और भारों को दो
तरह से समूहबद्ध कीजिए:
\[
G = \operatorname{bar}\bigl(M_{AB}, 2;\ M_{CD}, 2\bigr)
= \operatorname{bar}\bigl(M_{BC}, 2;\ M_{DA}, 2\bigr) :
\]
अर्थात् $G$ \emph{दोनों} द्वि-माध्यिकाओं का मध्यबिंदु है --- इसलिए दोनों
द्वि-माध्यिकाएँ सदा एक-दूसरे को समद्विभाजित करती हैं, और चारों मध्यबिंदुओं
का चतुर्भुज समांतर चतुर्भुज है (उसके विकर्ण वही द्वि-माध्यिकाएँ हैं)। न
कोई स्थिति-विश्लेषण, न निर्देशांक, और यह तर्क $\R^3$ में किसी विषमतलीय चतुर्भुज
के लिए भी ज्यों का त्यों चलता है, जहाँ चित्र पर आधारित उपपत्ति पहले ही
नाज़ुक हो जाती: साहचर्यता को विमा से कोई मतलब नहीं।
\end{example}

\begin{example}[किसी आनत प्रतिचित्रण का वर्गीकरण, आरंभ से अंत तक]\label{ex:b2:affine:dilation}
$\R^2$ पर $f(x, y) = (2x - 1,\ 3y - 4)$ लीजिए। उसका रैखिक भाग $\varphi = \operatorname{diag}(2, 3)$ है, जिसका \omterm{def:b2:reduction:eigen}{वर्णक्रम} $\{2, 3\}$ $1$ से बचा
रहता है: अतः नीचे सिद्ध की गई अचल-बिंदु कसौटी (\cref{prop:b2:affine:fixedpoint}) से $f$ का ठीक एक अचल
बिंदु है, जो
\[
x = 2x - 1, \qquad y = 3y - 4
\qquad\Longrightarrow\qquad \Omega = (1, 2).
\]
हल करने पर मिलता है। $\Omega$ पर पुनःकेंद्रण ($x = 1 + u$, $y = 2 + v$ रखिए):
\[
f(1 + u,\ 2 + v) = (1 + 2u,\ 2 + 3v) :
\]
अर्थात् $\Omega$ वाले ढाँचे में $f$ \emph{ही} उसका रैखिक भाग है: एक असमदैशिक
विस्फारण, जो केंद्र $(1, 2)$ से क्षैतिज दिशा में $2$ और ऊर्ध्वाधर में $3$
गुना खींचता है। सामान्य पाठ: कोई \omterm{def:b2:affine:subspace}{आनत प्रतिचित्रण} ``रैखिक प्रतिचित्रण जमा
स्थान-सूचना'' है, और जब भी $1$ \omterm{def:b2:reduction:eigen}{अभिलक्षणिक मान} न हो तब वह स्थान-सूचना किसी
एक अच्छी तरह चुने गए मूल बिंदु पर सिमट जाती है। विलोमतः मूल बिंदु को ग़लत
जगह ले जाने से अचर पद \emph{बन} जाते हैं: \omterm{def:b2:affine:subspace}{आनत} ज्यामिति यह चुनने की कला है
कि $0$ कहाँ रखा जाए।
\end{example}

\begin{remark}[विधि: बैरिकेंद्रों से संगामिता और संरेखता]
\cref{ex:b2:affine:median} इसी सामान्य विधि का एक उदाहरण है। किसी त्रिभुज की तीन चेवियन संगामी
हैं, यह सिद्ध करने के लिए कोई \emph{एक} भारित निकाय $(A,
\alpha; B, \beta; C, \gamma)$ दिखाइए और
साहचर्यता तीन तरह से लगाइए: $(B, C)$ समूहबद्ध करने पर \omterm{def:b2:affine:barycenter}{बैरिकेंद्र} $A$ से आने
वाली चेवियन पर दिखता है, $(C, A)$ समूहबद्ध करने पर $B$ से आने वाली पर, और $(A, B)$
समूहबद्ध करने पर तीसरी पर। माध्यिकाओं के लिए निकाय $(A, 1; B, 1; C, 1)$ सारा काम कर देता
है; और नियत अनुपातों में भुजाएँ काटने वाली चेवियनों के लिए भार उन्हीं
अनुपातों से पढ़ लिए जाते हैं। और तीन बिंदुओं की \emph{संरेखता} सिद्ध करने
के लिए एक को शेष दो का \omterm{def:b2:affine:barycenter}{बैरिकेंद्र} लिखिए (\cref{exo:b2:affine:2}), अथवा \cref{exo:b2:affine:11} की सारणिक-कसौटी
काम में लीजिए। दोनों विधियाँ ज्यामितीय चतुराई के स्थान पर भार का लेखा-जोखा
रख देती हैं --- और बैरिकेंद्रीय कलन ठीक इसी के लिए है।
\end{remark}

\section{उत्तलता, आनत रूप में}

\begin{definition}\label{def:b2:affine:convex}
किसी \omterm{def:b2:affine:def}{आनत समष्टि} का उपसमुच्चय $C$ \emph{उत्तल}\index{उत्तल समुच्चय} कहलाता
है जब वह अपने बिंदुओं के \emph{अऋणात्मक} भारों वाले हर \omterm{def:b2:affine:barycenter}{बैरिकेंद्र} को
समाहित करे --- समतुल्य रूप से, अपने बिंदुओं के बीच का हर खंड $\intcc{A}{B} = \{\operatorname{bar}(A,
1-t; B, t) : t \in \intcc{0}{1}\}$।
\emph{उत्तल कवच}\index{उत्तल कवच} $\operatorname{conv}(S)$ $S$ के बिंदुओं के सभी अऋणात्मक-भार
बैरिकेंद्रों का समुच्चय है --- अर्थात् $S$ को समाहित करने वाला सबसे छोटा
उत्तल समुच्चय।
\end{definition}

\begin{example}[अधिआलेख उत्तल समुच्चय होते हैं]\label{ex:b2:affine:epigraph}
परवलय के ऊपर का क्षेत्र $C = \{(x, y) : y \geq x^2\}$ \omterm{def:b2:affine:convex}{उत्तल} है: $(x_1, y_1), (x_2, y_2) \in C$ और $t \in
\intcc01$ के लिए वर्ग फलन की
उत्तलता असमिका देती है
\[
\bigl((1-t)x_1 + tx_2\bigr)^2 \leq (1-t)x_1^2 + tx_2^2
\leq (1-t)y_1 + ty_2 ,
\]
अतः \omterm{def:b2:affine:barycenter}{बैरिकेंद्र} परवलय के ऊपर ही रहता है। यह परिकलन सामान्य है: $\{y \geq f(x)\}$ ठीक तब
\omterm{def:b2:affine:convex}{उत्तल} है जब $f$ उत्तल फलन हो --- अर्थात् \omterm{def:b2:affine:convex}{उत्तल} \emph{समुच्चय} और \omterm{def:b2:affine:convex}{उत्तल}
\emph{फलन} (\cref{ch:b2:realfun}) एक ही धारणा के दो चेहरे हैं, और अधिआलेख उनका शब्दकोश। और
यही वह ज्यामितीय कारण है कि उत्तल फलनों की अवलंब रेखाएँ विद्यमान होती हैं
--- वही तथ्य जो \cref{ch:b2:randomvar} में जेनसन की असमिका सिद्ध करेगा।
\end{example}

\begin{example}[किसी उत्तल कवच के अनावश्यक जनक]\label{ex:b2:affine:squarehull}
$S = \{(0,0), (2,0), (2,2), (0,2), (1,1)\}$ लीजिए। पाँचवाँ बिंदु \omterm{def:b2:affine:barycenter}{बैरिकेंद्र}
\[
(1,1) = \operatorname{bar}\bigl((0,0), \tfrac12;\ (2,2),
\tfrac12\bigr),
\]
है, अतः वह पहले ही शेष चारों के कवच में पड़ा है: $\operatorname{conv}(S)$ वही वर्ग है जिसके
चारों कोने शीर्ष हैं। सामान्य रूप में $S$ का वह बिंदु जो $S$ के
\emph{शेष} बिंदुओं का अऋणात्मक-भार \omterm{def:b2:affine:barycenter}{बैरिकेंद्र} हो, कवच बदले बिना हटाया जा
सकता है; और जिन बिंदुओं को कभी नहीं हटाया जा सकता (यहाँ चारों कोने) वे कवच
के \emph{चरम बिंदु} हैं। उनका निर्धारण विशुद्ध बैरिकेंद्रीय परिकलन है: मान
लीजिए $(2,0)$ को शेष बिंदुओं के अऋणात्मक भारों वाले $\operatorname{bar}$ के रूप में नहीं लिखा
जा सकता, क्योंकि पहला निर्देशांक सारा भार $x = 2$ वाले बिंदुओं पर डाल देता है
और तब दूसरा निर्देशांक विफल हो जाता है। उत्तलता के प्रश्न बार-बार छोटे
भारित निकाय हल करने पर सिमट आते हैं।
\end{example}

\begin{theorem}[कारेथेओदोरी]\label{thm:b2:affine:caratheodory}
विमा $n$ की \omterm{def:b2:affine:def}{आनत समष्टि} में $\operatorname{conv}(S)$ का हर बिंदु $S$ के अधिकतम $n + 1$ बिंदुओं का
\omterm{def:b2:affine:barycenter}{बैरिकेंद्र} होता है।\index{कारेथेओदोरी की प्रमेय}
\end{theorem}

\begin{proof}
मान लीजिए $\lambda_i > 0$, $\sum\lambda_i = 1$ तथा $k + 1 > n + 1$ बिंदुओं के साथ $G = \operatorname{bar}(A_0, \lambda_0; \dots; A_k, \lambda_k)$। $k$ सदिश $\vect{A_0A_i}$ ($i \geq 1$) परस्पर
जुड़े हैं ($k > n$): अतः $\sum_{i\geq1}\mu_i \vect{A_0A_i} = 0$ अतुच्छ रूप से; और $\mu_0 = -\sum_{i\geq1}\mu_i$ रखने पर भार $(\mu_i)$ मिलते हैं
जिनके लिए $\sum\mu_i = 0$, $\sum \mu_i\,\vect{OA_i} = 0$ (कोई भी $O$), और सब शून्य नहीं। तब प्रत्येक वास्तविक
$t$ के लिए भार $\lambda_i - t\mu_i$ का योग तब भी $1$ रहता है, और चूँकि $\sum_i\mu_i\vect{OA_i} = 0$,
\[
\sum_i(\lambda_i - t\mu_i)\,\vect{OA_i}
= \sum_i\lambda_i\,\vect{OA_i} :
\]
वे \emph{वही} बिंदु $G$ देते हैं। अब $t$ को $0$ से सरकाइए: कोई $\mu_i$
धनात्मक है (उनका योग शून्य है और सब शून्य नहीं), अतः
\[
t^* = \min\Bigl\{\frac{\lambda_i}{\mu_i} : \mu_i > 0\Bigr\}
\]
सुपरिभाषित और धनात्मक है। $t = t^*$ पर: $\mu_i > 0$ वाले सूचकांकों के लिए न्यूनतमता से
$\lambda_i - t^*\mu_i \geq 0$, और किसी न्यूनतमकारी सूचकांक पर समता; और $\mu_i
\leq 0$ वाले सूचकांकों के लिए
$\lambda_i - t^*\mu_i \geq \lambda_i > 0$। इस प्रकार सारे भार अऋणात्मक बने रहते हैं और कम से कम एक मर जाता है:
अर्थात् $G$ कम बिंदुओं के \omterm{def:b2:affine:barycenter}{बैरिकेंद्र} के रूप में फिर से लिखा जाता है। और
जब तक $n + 1$ से अधिक बिंदु बचे हों तब तक दोहराइए।
\end{proof}

\begin{example}
समतल में ($n = 2$): किसी परिमित समुच्चय के \omterm{def:b2:affine:convex}{उत्तल कवच} का हर बिंदु उसी समुच्चय
के शीर्षों वाले किसी त्रिभुज में पड़ता है --- यही कारेथेओदोरी की ज्यामितीय
सामग्री है, जो अनुकूलन और प्रायिकता (मिश्रण) दोनों में काम आती है।
\end{example}

\begin{example}[कारेथेओदोरी की कलनविधि चलाकर]\label{ex:b2:affine:caratheodoryrun}
\cref{ex:b2:affine:squarehull} के वर्ग का केंद्र उसके चारों कोनों $A_1 =
(0,0)$, $A_2 = (2,0)$, $A_3 = (2,2)$, $A_4 = (0,2)$ से लिखिए:
\[
(1,1) = \operatorname{bar}\bigl(A_1, \tfrac14;\ A_2,
\tfrac14;\ A_3, \tfrac14;\ A_4, \tfrac14\bigr),
\]
अर्थात् विमा $2$ में चार बिंदु --- एक अधिक। उपपत्ति की विधि $\sum\mu_i = 0$ तथा $\sum\mu_i\vect{OA_i} = 0$
वाले भार $(\mu_i)$ माँगती है: यहाँ $\mu = (1, -1, 1, -1)$ काम कर जाता है (दोनों विकर्णों का
मध्यबिंदु साझा है)। $\lambda_i
\mapsto \lambda_i - t\mu_i$ को सरकाने पर प्रत्येक $t$ के लिए \omterm{def:b2:affine:barycenter}{बैरिकेंद्र}
अपरिवर्तित रहता है; और चरम स्वीकार्य मान $t = \frac14$ भारों को $(0, \tfrac12, 0, \tfrac12)$ बना देता है,
जो $A_1$ और $A_3$ दोनों को एक साथ मार देता है:
\[
(1,1) = \operatorname{bar}\bigl(A_2, \tfrac12;\ A_4,
\tfrac12\bigr),
\]
अर्थात् दो बिंदुओं से निरूपण --- और यह प्रमेय की गारंटी वाले तीन से भी
बेहतर है, क्योंकि यहाँ केंद्र संयोग से जनकों के बीच के किसी खंड पर पड़ता
है। कलनविधि पूरी तरह यांत्रिक है: कोई परतंत्रता ढूँढ़िए, तब तक सरकाइए जब
तक कोई भार मर जाए, और दोहराइए।
\end{example}

\section{आनत वर्गीकरण के औज़ार}

\begin{proposition}[आनत प्रतिचित्रणों के अचल बिंदु]\label{prop:b2:affine:fixedpoint}
मान लीजिए $f$ किसी परिमित-विमीय \omterm{def:b2:affine:def}{आनत समष्टि} का \omterm{def:b2:affine:subspace}{आनत} अंतःरूपांतरण है जिसका
रैखिक भाग $\varphi$ है। यदि $1 \notin
\operatorname{Sp}(\varphi)$, तो $f$ का ठीक एक अचल बिंदु $\Omega$ है, और $\Omega$ पर
सदिशीकरण में $f$ \emph{ही} उसका रैखिक भाग है। (स्थानांतरण, जिनके लिए $\varphi = \mathrm{id}$
और कोई अचल बिंदु नहीं, मूल बाधा हैं।)
\end{proposition}

\begin{proof}
$O$ स्थिर कीजिए और $f(O + x) = f(O) + \varphi(x)$ लिखिए। बिंदु $O + x$ अचल है तभी जब $O + x = f(O) + \varphi(x)$, अर्थात्
\[
(\mathrm{id} - \varphi)(x) = \vect{O f(O)} .
\]
परिमित विमा में $\mathrm{id} - \varphi$ प्रतिलोमनीय है तभी जब $0$ $\mathrm{id} - \varphi$ का \omterm{def:b2:reduction:eigen}{अभिलक्षणिक मान} न हो,
अर्थात् तभी जब $1 \notin \operatorname{Sp}\varphi$ --- और उस स्थिति में प्रदर्शित समीकरण का ठीक एक हल $x^*$
है, जो अद्वितीय अचल बिंदु $\Omega = O + x^*$ देता है। पुनःकेंद्रण: किसी भी सदिश $u$ के
लिए
\[
f(\Omega + u) = f(\Omega) + \varphi(u) = \Omega + \varphi(u),
\]
अतः मूल बिंदु $\Omega$ वाले ढाँचे में प्रतिचित्रण $u \mapsto
\varphi(u)$ पढ़ा जाता है: अर्थात्
विशुद्ध रैखिक। और $1 \in
\operatorname{Sp}\varphi$ होने पर या तो कोई अचल बिंदु विद्यमान ही नहीं होता
(प्रदर्शित समीकरण अहल्य हो सकता है, जैसे किसी स्थानांतरण के लिए) या उनकी
पूरी \omterm{def:b2:affine:subspace}{आनत} उपसमष्टि होती है (किसी हल में \omterm{def:b2:reduction:eigen}{अभिलक्षणिक मान} $1$ का कोई भी
\omterm{def:b2:reduction:eigen}{अभिलक्षणिक सदिश} जोड़ दीजिए): अतः अद्वितीयता ठीक वही वर्णक्रमीय शर्त है।
\end{proof}

\begin{example}[समतलीय समदूरिक प्रतिचित्रण, पूरे]\label{ex:b2:affine:isometries}
यूक्लिडीय समतल के किसी \omterm{def:b2:affine:subspace}{आनत} समदूरिक प्रतिचित्रण का रैखिक भाग $O(2)$ में होता
है: कोई घूर्णन $R_\theta$ अथवा कोई परावर्तन (प्रथम वर्ष का खंड)। यदि $\theta \neq 0$: तो $1 \notin \operatorname{Sp} R_\theta$,
अतः प्रतिचित्रण \emph{किसी अद्वितीय केंद्र के परितः घूर्णन} है (\cref{prop:b2:affine:fixedpoint})। और
यदि रैखिक भाग कोई परावर्तन हो: तो या तो किसी अक्ष में परावर्तन (अचल बिंदु
विद्यमान) अथवा कोई \emph{सर्पी परावर्तन} (अक्ष के अनुदिश स्थानांतरण के साथ
संयोजित परावर्तन, बिना किसी अचल बिंदु के)। स्थानांतरणों को मिलाकर यही
समतलीय समदूरिक प्रतिचित्रणों का पूरा वर्गीकरण है।
\end{example}

\begin{remark}[समतलीय समदूरिक प्रतिचित्रण, एक नज़र में]
स्थितियाँ इकट्ठी कीजिए: तत्समक; स्थानांतरण ($\vec f =
\mathrm{id}$, और तुच्छ को छोड़कर कोई
अचल बिंदु नहीं); घूर्णन (रैखिक भाग $R_\theta$, $\theta \neq 0$: एक केंद्र); परावर्तन (रैखिक
भाग कोई परावर्तन, और अचल बिंदुओं की एक रेखा); तथा सर्पी परावर्तन (वही
रैखिक भाग, पर कोई अचल बिंदु नहीं)। अर्थात् तत्समक के अतिरिक्त चार कुल, और
हर एक केवल दो आँकड़ों से पहचाना जाता है: रैखिक भाग और अचल-बिंदु समुच्चय ---
यानी \cref{prop:b2:affine:fixedpoint} का ढर्रा, पूरी तरह गिनाया हुआ।
\end{remark}

\begin{example}[एक सर्पी परावर्तन, रंगे हाथ पकड़ा हुआ]\label{ex:b2:affine:glide}
$f(x, y) = (y + 1,\ x + 1)$ लीजिए। रैखिक भाग $(x, y)
\mapsto (y, x)$ विकर्ण $y = x$ में परावर्तन है, अतः $1 \in \operatorname{Sp}\vec f$ और \cref{prop:b2:affine:fixedpoint} चुप
है। अचल बिंदुओं के लिए एक साथ $x = y + 1$ तथा $y = x + 1$ चाहिए: जो असंभव है --- अतः कोई
अचल बिंदु नहीं, इसलिए $f$ परावर्तन नहीं है। वर्ग करने से वर्गीकरण तय हो
जाता है:
\[
f\bigl(f(x, y)\bigr) = f(y + 1,\ x + 1) = (x + 2,\ y + 2),
\]
अर्थात् $(2, 2)$ से स्थानांतरण: इसलिए $f$ वह \emph{सर्पी परावर्तन} है जिसका
अक्ष रेखा $y = x$ है (उपयुक्त रूप से खिसकाई हुई: $M$ और $f(M)$ का मध्यबिंदु
सदा $y = x + {}$अचर पर पड़ता है, यहाँ $y = x$, जैसा $M = (0,
0) \mapsto (1,1)$ पर जाँचा जा सकता है) और
सर्पण सदिश $(1, 1)$ है, जो $f
\circ f$ का आधा है। इसकी तुलना \cref{exo:b2:affine:6} से कीजिए, जहाँ वही
रैखिक भाग पर भिन्न अचर पद सच्चा परावर्तन देता है: \omterm{def:b2:reduction:eigen}{अभिलक्षणिक मान} $1$ के
उपस्थित होने पर अचर पद ही सब कुछ तय करता है।
\end{example}

\begin{example}[आनत पुनरावृत्तियाँ आनत गतिकी हैं]\label{ex:b2:affine:recursion}
क्लासिक पुनरावृत्ति $u_{n+1} = au_n + b$ ($a \neq 1$) रेखा के \omterm{def:b2:affine:subspace}{आनत प्रतिचित्रण} $f(x) = ax + b$ को दोहराती है,
जिसका रैखिक भाग $a$ \omterm{def:b2:reduction:eigen}{अभिलक्षणिक मान} $1$ से बचा रहता है: अतः कोई अद्वितीय
अचल बिंदु $\omega = \frac{b}{1-a}$ है, और वहाँ पुनःकेंद्रण करने पर (ऊपर वाली प्रतिज्ञप्ति की
एक-विमीय स्थिति) $f$ $a$ से गुणन बन जाता है:
\[
u_{n+1} - \omega = a\,(u_n - \omega)
\qquad\Longrightarrow\qquad
u_n = \omega + a^n(u_0 - \omega) .
\]
$u_{n+1} = \frac{u_n}2 + 3$ के लिए: $\omega = 6$ और $u_n = 6 +
(u_0 - 6)2^{-n} \to 6$। ऐसी पुनरावृत्तियों के लिए \cref{ch:b2:series} में सिखाई गई
विधि --- ``अचल बिंदु घटा दीजिए'' --- ठीक किसी \omterm{def:b2:affine:subspace}{आनत प्रतिचित्रण} का उसके अचल
बिंदु पर सदिशीकरण है; और $\abs a < 1$ के लिए अभिसरण वही संकुचन-परिघटना है जिसे \cref{ch:b2:metric}
ने बानाख अचल बिंदु प्रमेय बना दिया। एक विचार, तीन अध्याय।
\end{example}

\begin{example}[किसी घूर्णन का केंद्र ढूँढ़ना]\label{ex:b2:affine:rotationcenter}
$f(x, y) = (-y + 2,\ x)$ लीजिए। रैखिक भाग $\varphi(x, y) = (-y, x)$ है: अर्थात् कोण $\frac\pi2$ का घूर्णन, जिसका \omterm{def:b2:reduction:eigen}{वर्णक्रम} $\{\iu, -\iu\}$
$1$ से बचा रहता है। \cref{prop:b2:affine:fixedpoint} के अनुसार ठीक एक अचल बिंदु है: $x = -y + 2$ और $y = x$ $x = 1$,
$y = 1$ देते हैं, अतः $\Omega = (1, 1)$, और $f$ केंद्र $(1, 1)$ तथा कोण $\frac\pi2$ का घूर्णन
\emph{ही} है। सामान्य पाठ: जब $1
\notin \operatorname{Sp}\vec f$ हो, तब $f$ के वर्गीकरण की क़ीमत एक
रैखिक निकाय है --- ज्यामिति पूरी तरह रैखिक भाग में है और अंकगणित पूरी तरह
केंद्र ढूँढ़ने में।
\end{example}

\begin{remark}[आनत भाषा आगे कहाँ काम आती है]
\omterm{def:b2:affine:barycenter}{बैरिकेंद्र} और \omterm{def:b2:affine:subspace}{आनत प्रतिचित्रण} आगे के ज्यामिति अध्यायों का व्याकरण हैं:
स्पर्श रेखाएँ और समतल \omterm{def:b2:affine:subspace}{आनत} वस्तुएँ हैं (\cref{ch:b2:curves,ch:b2:surfaces}), \omterm{def:b2:affine:subspace}{आनत} चर-परिवर्तन क्षेत्रफल और
आयतन को $\abs{\det \vec f}$ से गुणा कर देता है (\cref{ch:b2:multint}), और प्रत्याशा किसी प्रायिकता वितरण
से मिले भारों वाला \omterm{def:b2:affine:barycenter}{बैरिकेंद्र} है, और इसीलिए उत्तलता जेनसन की असमिका पर
शासन करती है (\cref{ch:b2:randomvar})। तृतीय वर्ष के खंड में यही उत्तलता-शब्दावली $L^p$
मानदंडों तथा समाकल असमिकाओं का अध्ययन उठाए चलती है।
\end{remark}

\begin{remark}[इसी खंड के भीतर के परिप्रेक्ष्य]
इस अध्याय से दो धागे निकलते हैं। \emph{\omterm{def:b2:affine:subspace}{आनत}} धागा: स्पर्श रेखाएँ (\cref{ch:b2:curves}) और
स्पर्श समतल (\cref{ch:b2:surfaces}) अरैखिक वस्तुओं से जुड़ी \omterm{def:b2:affine:subspace}{आनत} उपसमष्टियाँ हैं, और पृष्ठ
वाले अध्याय में द्विघातीय पृष्ठों का वर्गीकरण इसी अध्याय के केंद्र-समीकरण
$A\Omega = -b$ पर चलता है। और \emph{\omterm{def:b2:affine:convex}{उत्तल}} धागा लंबा है: अर्धसमतलों तथा चक्रिकाओं की
उत्तलता सप्ताहांत समस्या का हेली-सिद्धांत चलाती है; फलनों की उत्तलता जेनसन
की असमिका देती है (\cref{ch:b2:randomvar}); और पुस्तक की अंतिम प्रमेय --- शाखन प्रक्रियाओं
के लिए विलोपन-कसौटी (\cref{ch:b2:genfun}) --- विकर्ण के सापेक्ष किसी उत्तल वक्र की स्थिति
से तय होती है, और वह चित्र जितना प्रायिकता का है उतना ही इस अध्याय का भी।
\omterm{def:b2:affine:barycenter}{बैरिकेंद्र} भी वहाँ लौटते हैं: कोई प्रत्याशा प्रायिकता-भारों वाला \omterm{def:b2:affine:barycenter}{बैरिकेंद्र}
ही है।
\end{remark}

\section{अभ्यास}

\begin{exercise}[$\star$]\label{exo:b2:affine:1}
$\R^3$ में क्या निम्नलिखित \omterm{def:b2:affine:subspace}{आनत} उपसमष्टियाँ हैं? दिशाएँ और विमाएँ दीजिए। $\{x + y + z = 1\}$;
$\;\{x + y + z = 1,\ x - z =
3\}$; $\;\{x^2 + y^2 = 1\}$; और किसी दिए हुए संगत निकाय $MX = B$ का हल-समुच्चय।
\end{exercise}

\begin{exercise}[$\star$]\label{exo:b2:affine:2}
सिद्ध कीजिए कि किसी \omterm{def:b2:affine:def}{आनत समष्टि} के तीन भिन्न बिंदु $A, B, C$ संरेख हैं तभी जब $C$
$A$ और $B$ का \omterm{def:b2:affine:barycenter}{बैरिकेंद्र} हो, तभी जब सदिश $\vect{AB}, \vect{AC}$ परस्पर जुड़े हों। इससे
मेनेलौस-शैली का भार-लेखा निष्कर्ष रूप में निकालिए: यदि $C = \operatorname{bar}(A, 1 - t; B, t)$, तो $t = \frac12$, $t = 2$,
$t = -1$ के लिए $C$ ज्ञात कीजिए।
\end{exercise}

\begin{exercise}[$\star$]\label{exo:b2:affine:3}
मान लीजिए $f$ $\R^2$ का वह \omterm{def:b2:affine:subspace}{आनत प्रतिचित्रण} है जो $M = \frac12\begin{pmatrix} 1 & 1\\ 1 & 1\end{pmatrix}$ तथा $C =
(1, 0)^{\mathsf T}$ के साथ $f(X) = MX + C$ से
दिया गया है। $f$ का प्रतिबिंब, उसके अचल बिंदु (यदि कोई हों), और $f \circ f$
निर्धारित कीजिए।
\end{exercise}

\begin{exercise}[$\star\star$]\label{exo:b2:affine:4}
(साहचर्यता क्रिया में) किसी त्रिभुज $ABC$ में मान लीजिए $I, J, K$ $BC$, $CA$, $AB$
को क्रमशः अनुपातों $\vect{BI} = \frac13\vect{BC}$, $\vect{CJ} = \frac13\vect{CA}$, $\vect{AK} = \frac13\vect{AB}$ में बाँटते हैं। $I, J, K$ को बैरिकेंद्रों के
रूप में व्यक्त कीजिए और $(I,1;J,1;K,1)$ का \omterm{def:b2:affine:barycenter}{बैरिकेंद्र} परिकलित कीजिए: आपको क्या मिलता
है, और वह पहले से ही अनुमेय क्यों था?
\end{exercise}

\begin{exercise}[$\star\star$]\label{exo:b2:affine:5}
सिद्ध कीजिए कि \emph{मध्यबिंदु} सुरक्षित रखने वाला ($f\bigl(\frac{A+B}{2}\bigr) = \frac{f(A) + f(B)}{2}$) और \emph{\omterm{def:b2:metric:continuity}{संतत}}
प्रतिचित्रण $f \colon \R^n \to \R^n$ \omterm{def:b2:affine:subspace}{आनत} होता है। \emph{(दिखाइए कि सदिश प्रतिचित्रण $u \mapsto
f(O + u) - f(O)$
मध्यबिंदुओं के द्वारा योज्य है, फिर $\Q$-समघातीय, फिर संततता से
$\R$-समघातीय --- वही सघनता वाली रणनीति जो प्रथम वर्ष के खंड में कोशी के
फलनात्मक समीकरण के लिए प्रयुक्त हुई थी; आवश्यक पग यहीं फिर से निकालिए।)}
\end{exercise}

\begin{exercise}[$\star\star$]\label{exo:b2:affine:6}
यूक्लिडीय समतल के \omterm{def:b2:affine:subspace}{आनत प्रतिचित्रण} $f(x, y) = (y + 1,\; x - 1)$ का वर्गीकरण कीजिए: रैखिक भाग, अचल
बिंदु, ज्यामितीय प्रकृति (परावर्तन? सर्पी?)। $f \circ f$ परिकलित कीजिए और निष्कर्ष
निकालिए।
\end{exercise}

\begin{exercise}[$\star\star\star$]\label{exo:b2:affine:7}
(रादों) मान लीजिए $A_1, \dots, A_{n+2}$ विमा $n$ की किसी \omterm{def:b2:affine:def}{आनत समष्टि} के बिंदु हैं। सिद्ध
कीजिए कि उन्हें ऐसे दो असंयुक्त समूहों में बाँटा जा सकता है जिनके \omterm{def:b2:affine:convex}{उत्तल
कवच} प्रतिच्छेद करते हैं। \emph{(कारेथेओदोरी की उपपत्ति की तरह ऐसे भार $\mu_i$
ढूँढ़िए जो सब शून्य न हों और $\sum\mu_i = 0$ तथा $\sum\mu_i\vect{OA_i} = 0$ पूरा करें; फिर धनात्मक और ऋणात्मक
भार अलग कीजिए और दोनों पक्षों का सामान्यीकरण कीजिए।)}
\end{exercise}

\begin{exercise}[$\star\star\star$]\label{exo:b2:affine:8}
मान लीजिए $f$ $\R^n$ का \omterm{def:b2:affine:subspace}{आनत} अंतःरूपांतरण है जहाँ $f \circ f = f$। सिद्ध कीजिए कि $f$
दिशा $\ker\vec f$ के अनुदिश \omterm{def:b2:affine:subspace}{आनत} उपसमष्टि $\operatorname{Fix}(f) = \operatorname{im} f$ पर \omterm{def:b2:affine:subspace}{आनत} प्रक्षेप है, और विलोमतः ऐसे
सारे प्रक्षेप वर्गसम होते हैं। \emph{(पहले दिखाइए कि $\operatorname{im} f$ अचल बिंदुओं से
बना है।)}
\end{exercise}

\begin{exercise}[$\star$]\label{exo:b2:affine:9}
किसी त्रिभुज $ABC$ में $G = \operatorname{bar}(A, 1;\ B, 2;\ C, 3)$ लीजिए। साहचर्यता का उपयोग करके दिखाइए कि रेखा $AG$
$BC$ से $M = \operatorname{bar}(B, 2;\ C, 3)$ पर मिलती है, और $G$ को खंड $\intcc AM$ पर ज्ञात कीजिए; इसी प्रकार $BG$
तथा $CA$ का प्रतिच्छेद भी ज्ञात कीजिए।
\end{exercise}

\begin{exercise}[$\star\star$]\label{exo:b2:affine:10}
$\lambda \neq 0$ के लिए \emph{समरूपण} $h_{\Omega,
\lambda}$ वह \omterm{def:b2:affine:subspace}{आनत प्रतिचित्रण} है जो $\Omega$ को अचल रखता है
और जिसका रैखिक भाग $\lambda\,\mathrm{id}$ है। सिद्ध कीजिए कि संयोजन $h_{\Omega', \mu} \circ h_{\Omega, \lambda}$ $\lambda\mu \neq 1$ होने पर अनुपात
$\lambda\mu$ का समरूपण है, और $\lambda\mu = 1$ होने पर स्थानांतरण; तथा $\lambda = \mu = -1$ वाली स्थिति (दो बिंदु
परावर्तन) में स्थानांतरण सदिश परिकलित कीजिए।
\end{exercise}

\begin{exercise}[$\star\star$]\label{exo:b2:affine:11}
(मेनेलौस) किसी त्रिभुज $ABC$ में $A' \in (BC)$, $B' \in
(CA)$, $C' \in (AB)$ लीजिए, जो सब शीर्षों से भिन्न
हों, और $\vect{A'B} =
\alpha\,\vect{A'C}$, $\vect{B'C} = \beta\,\vect{B'A}$, $\vect{C'A} = \gamma\,\vect{C'B}$ से $\alpha, \beta, \gamma$ परिभाषित कीजिए। सिद्ध कीजिए कि $A', B', C'$ संरेख हैं
यदि और केवल यदि $\alpha\beta\gamma = 1$। \emph{(हर बिंदु को दो शीर्षों का \omterm{def:b2:affine:barycenter}{बैरिकेंद्र} लिखिए;
और दिखाइए कि तीन बिंदु संरेख हैं तभी जब $(A, B, C)$ के सापेक्ष उनकी बैरिकेंद्रीय
निर्देशांक-पंक्तियाँ कोई अपभ्रष्ट $3 \times 3$ आव्यूह बनाएँ।)}
\end{exercise}

\begin{exercise}[$\star\star\star$]\label{exo:b2:affine:12}
सिद्ध कीजिए कि $\R^n$ के किसी \omterm{def:b2:metric:compact}{संहत} उपसमुच्चय $K$ का \omterm{def:b2:affine:convex}{उत्तल कवच} \omterm{def:b2:metric:compact}{संहत} होता है।
\emph{(\cref{thm:b2:affine:caratheodory} के अनुसार $\operatorname{conv}(K)$ किसी \omterm{def:b2:metric:compact}{संहत} समुच्चय का \omterm{def:b2:metric:continuity}{संतत} प्रतिचित्रण के अंतर्गत
प्रतिबिंब है।)} $\R^2$ में किसी उदाहरण से दिखाइए कि किसी \emph{संवृत} समुच्चय
का \omterm{def:b2:affine:convex}{उत्तल कवच} संवृत होना आवश्यक नहीं।
\end{exercise}

\section{समस्या: रादों से हेली तक, केंद्रबिंदु और युंग की प्रमेय}

\begin{omfigure}
\begin{tikzpicture}[scale=0.9]
  \draw[omDef, thick] (0,0) -- (3,0) -- (0.9,2.6) -- cycle;
  \fill (0,0) circle (1.6pt) node[below left] {\small $A_1$};
  \fill (3,0) circle (1.6pt) node[below right] {\small $A_2$};
  \fill (0.9,2.6) circle (1.6pt) node[above] {\small $A_3$};
  \fill[omThm] (1.3,0.9) circle (1.8pt)
    node[below, omThm] {\small $A_4$};
\end{tikzpicture}
\qquad
\begin{tikzpicture}[scale=0.9]
  \draw[omDef, thick] (0,0) -- (3,0.3) -- (3.3,2.4) --
    (0.4,2.2) -- cycle;
  \draw[omThm, thick] (0,0) -- (3.3,2.4);
  \draw[omThm, thick] (3,0.3) -- (0.4,2.2);
  \fill (0,0) circle (1.6pt) node[below left] {\small $A_1$};
  \fill (3,0.3) circle (1.6pt) node[below right] {\small $A_2$};
  \fill (3.3,2.4) circle (1.6pt) node[above right] {\small $A_3$};
  \fill (0.4,2.2) circle (1.6pt) node[above left] {\small $A_4$};
  \fill[omProp] (1.75,1.27) circle (1.8pt);
\end{tikzpicture}

{\small व्यापक स्थिति में समतल के चार बिंदुओं के लिए दोनों रादों प्रकार:
या तो एक बिंदु शेष तीन के त्रिभुज के भीतर (विभाजन $\{A_4\} \mid \{A_1, A_2, A_3\}$), या \omterm{def:b2:affine:convex}{उत्तल} स्थिति,
जहाँ रादों बिंदु (नारंगी) दोनों विकर्णों का प्रतिच्छेद है।}
\end{omfigure}

\begin{problem}[{सप्ताहांत समस्या --- हेली की प्रमेय और उसके दो
लाभ}]\label{pb:b2:affine:1}
रादों की प्रमेयिका (\cref{exo:b2:affine:7}) कहती है कि $n$-विमीय \omterm{def:b2:affine:def}{आनत समष्टि} के $n + 2$ बिंदु सदा
ऐसे दो समूहों में बँट जाते हैं जिनके \omterm{def:b2:affine:convex}{उत्तल कवच} प्रतिच्छेद करते हैं। यह
समस्या रैखिक बीजगणित के उसी एक तथ्य को संचयात्मक ज्यामिति की प्रमेयों की
शृंखला में बदल देती है: हेली की प्रतिच्छेदन प्रमेय, केंद्रबिंदु प्रमेय (एक
द्विविमीय माध्यिका), और युंग की आच्छादन प्रमेय। आगे सर्वत्र समतल $\R^2$ है
अपनी सामान्य यूक्लिडीय संरचना के साथ, और $\det$ विहित आधार में \omterm{def:b2:linalg:det}{सारणिक}
है।\index{हेली की प्रमेय}\index{रादों की प्रमेयिका}\index{युंग की
प्रमेय}\index{केंद्रबिंदु}

\medskip
\noindent\textbf{भाग I --- बैरिकेंद्रीय निर्देशांक।} बिंदु $A_0, \dots, A_k$ \emph{\omterm{def:b2:affine:subspace}{आनत}
रूप से स्वतंत्र} कहलाते हैं जब सदिश $\vect{A_0A_1}, \dots, \vect{A_0A_k}$ रैखिकतः स्वतंत्र हों।
\begin{enumerate}
  \item दिखाइए कि \omterm{def:b2:affine:subspace}{आनत} स्वतंत्रता आधार-बिंदु $A_0$ के चुनाव पर निर्भर नहीं
        करती, और यह इसके तुल्य है: जब भी $1$ तक जुड़ने वाले भारों के दो
        कुल $(A_0,
        \dots, A_k)$ का एक ही \omterm{def:b2:affine:barycenter}{बैरिकेंद्र} परिभाषित करें, तब भार भी एक ही हों।
  \item मान लीजिए $(A, B, C)$ समतल में \omterm{def:b2:affine:subspace}{आनत} रूप से स्वतंत्र हैं। दिखाइए कि हर बिंदु
        $M$ के लिए $\alpha + \beta +
        \gamma = 1$ तथा $M = \operatorname{bar}(A, \alpha; B,
        \beta; C, \gamma)$ वाला अद्वितीय त्रिक $(\alpha, \beta, \gamma)$ है --- अर्थात्
        उसके \emph{बैरिकेंद्रीय निर्देशांक}।
  \item \omterm{def:b2:linalg:det}{सारणिक} सूत्र सिद्ध कीजिए
        \[
        \alpha = \frac{\det(\vect{MB},
        \vect{MC})}{\det(\vect{AB}, \vect{AC})},
        \qquad
        \beta = \frac{\det(\vect{MC},
        \vect{MA})}{\det(\vect{AB}, \vect{AC})},
        \qquad
        \gamma = \frac{\det(\vect{MA},
        \vect{MB})}{\det(\vect{AB}, \vect{AC})} :
        \]
        अर्थात् बैरिकेंद्रीय निर्देशांक चिह्नित क्षेत्रफलों के अनुपात हैं।
  \item रेखाएँ $BC$, $CA$, $AB$ निर्देशांक-रेखाएँ $\{\alpha = 0\}$, $\{\beta = 0\}$, $\{\gamma = 0\}$ हैं।
        दिखाइए कि $M$ संवृत त्रिभुज $\operatorname{conv}\{A, B, C\}$ में पड़ता है तभी जब $\alpha, \beta,
        \gamma \geq 0$, और यह
        कि तीनों रेखाएँ समतल को ठीक सात क्षेत्रों में काटती हैं, जिनका
        वर्गीकरण $(\alpha, \beta, \gamma)$ के चिह्नों से होता है (और चिह्न-प्रतिरूप $(-, -,
        -)$ असंभव
        है)।
  \item मान लीजिए $u \colon \R^2 \to \R$ कोई \omterm{def:b2:affine:subspace}{आनत प्रतिचित्रण} है (एक \emph{\omterm{def:b2:affine:subspace}{आनत} फलनिक})।
        दिखाइए $u(M) = \alpha\,u(A) +
        \beta\,u(B) + \gamma\,u(C)$, कि किसी अनचर \omterm{def:b2:affine:subspace}{आनत} फलनिक के स्तर समुच्चय रेखाएँ हैं,
        कि हर रेखा इसी तरह मिलती है, और यह कि संवृत अर्धसमतल $\{u
        \geq c\}$ \omterm{def:b2:affine:convex}{उत्तल}
        हैं।
\end{enumerate}

\medskip
\noindent\textbf{भाग II --- रादों विभाजन, परिष्कृत।} $\R^n$ के $n + 2$ बिंदुओं
का कुल \emph{व्यापक स्थिति} में तब कहलाता है जब उनमें से कोई भी $n + 1$ \omterm{def:b2:affine:subspace}{आनत}
रूप से स्वतंत्र हों। $(A_1, \dots, A_{n+2})$ की \emph{\omterm{def:b2:affine:subspace}{आनत} परतंत्रता} ऐसा कुल $(\mu_i)$ है जिसके लिए
$\sum_i \mu_i = 0$ और किसी (अतः हर) मूल बिंदु $O$ के लिए $\sum_i
\mu_i\,\vect{OA_i} = 0$।
\begin{enumerate}[resume]
  \item चार बिंदुओं $A_1 = (0,0)$, $A_2 = (3,0)$, $A_3 = (0,3)$, $A_4 = (1,1)$ की कोई अशून्य \omterm{def:b2:affine:subspace}{आनत} परतंत्रता परिकलित
        कीजिए; रादों विभाजन और रादों बिंदु दीजिए।
  \item दिखाइए कि व्यापक स्थिति के बिंदुओं के लिए \omterm{def:b2:affine:subspace}{आनत} परतंत्रताओं की सदिश
        समष्टि की विमा ठीक $1$ है, और किसी अशून्य परतंत्रता का कोई भी
        गुणांक शून्य \emph{नहीं} होता।
  \item इससे निष्कर्ष निकालिए कि व्यापक स्थिति के $n + 2$ बिंदुओं का रादों
        विभाजन अद्वितीय है (दोनों खंडों की अदला-बदली को छोड़कर), जहाँ हर
        खंड उन सूचकांकों का समुच्चय है जिन पर $\mu_i$ का चिह्न एक ही रहता है।
  \item समतल के व्यापक स्थिति वाले चार बिंदुओं के लिए द्विभाजन दिखाइए: या
        तो विभाजन का प्रकार $(1, 3)$ है --- अर्थात् एक बिंदु शेष तीन के
        त्रिभुज के भीतर --- या प्रकार $(2, 2)$: चारों बिंदु \omterm{def:b2:affine:convex}{उत्तल} स्थिति में
        हैं और दोनों युग्मों को जोड़ने वाले खंड (विकर्ण) रादों बिंदु पर
        प्रतिच्छेद करते हैं।
  \item प्रश्न 6 को इकाई वर्ग $(0,0)$, $(1,0)$, $(1,1)$, $(0,1)$ के लिए कीजिए:
        परतंत्रता, विभाजन, रादों बिंदु।
\end{enumerate}

\medskip
\noindent\textbf{भाग III --- समतल में हेली की प्रमेय।}
\begin{enumerate}[resume]
  \item मान लीजिए $C_1, C_2, C_3, C_4$ $\R^2$ के \omterm{def:b2:affine:convex}{उत्तल} उपसमुच्चय हैं जिनमें से किन्हीं तीन का
        कोई साझा बिंदु है। $x_i \in
        \bigcap_{j \neq i} C_j$ चुनिए और $x_1, \dots, x_4$ पर रादों की प्रमेयिका लगाइए:
        दिखाइए कि रादों बिंदु चारों समुच्चयों में है। \emph{(प्रत्येक $k$
        के लिए जिस खंड में $x_k$ नहीं है वह $C_k$ के बिंदुओं से बना है।)}
  \item (हेली) मान लीजिए $C_1, \dots, C_m$ ($m \geq 3$) $\R^2$ के \omterm{def:b2:affine:convex}{उत्तल} उपसमुच्चय हैं जिनमें से
        किन्हीं तीन का प्रतिच्छेद अरिक्त है। $m$ पर आगमन से $\bigcap_{i=1}^m C_i \neq \emptyset$ सिद्ध
        कीजिए: $C_{m-1}$ और $C_m$ के स्थान पर $C_{m-1} \cap C_m$ रखिए और नए कुल के लिए
        परिकल्पना प्रश्न 11 से जाँचिए।
  \item तीन प्रतिउदाहरण, हर परिकल्पना के लिए एक: (क) किसी त्रिभुज की तीनों
        संवृत भुजाएँ जोड़े-जोड़े में प्रतिच्छेद करती हैं पर उनका कोई साझा
        बिंदु नहीं ($3$ को $2$ तक नहीं घटाया जा सकता); (ख) व्यापक
        स्थिति के चार बिंदुओं के लिए चारों समुच्चय $S_i = \{x_1, \dots, x_4\}
        \setminus \{x_i\}$ त्रिक-प्रतिच्छेदन
        परिकल्पना पूरी करते हैं पर निष्कर्ष नहीं (उत्तलता मायने रखती है);
        (ग) संवृत अर्धसमतल $H_k = \intco{k}{+\infty} \times
        \R$, $k \in \N$, जोड़े-जोड़े और तीन-तीन में
        प्रतिच्छेद करते हैं पर $\bigcap_k H_k = \emptyset$ (अनंत कुलों को संहतता चाहिए)।
  \item (\omterm{def:b2:metric:compact}{संहत} हेली) मान लीजिए $(K_i)_{i \in I}$ $\R^2$ के \omterm{def:b2:metric:compact}{संहत} \omterm{def:b2:affine:convex}{उत्तल} उपसमुच्चयों का कोई
        भी कुल है जिनमें से किन्हीं तीन का प्रतिच्छेद अरिक्त है। प्रश्न 12
        और बोरेल--लेबेग गुणधर्म (\cref{thm:b2:metric:borellebesgue}) का उपयोग करके दिखाइए $\bigcap_{i \in I} K_i \neq \emptyset$।
  \item (पहला लाभ) मान लीजिए $S$ समतल के बिंदुओं का कोई परिमित समुच्चय है
        और $r > 0$। दिखाइए: यदि $S$ के किन्हीं तीन बिंदु त्रिज्या $r$ की
        किसी संवृत चक्रिका में पड़ते हों, तो $S$ पूरा त्रिज्या $r$ की
        किसी एक संवृत चक्रिका में पड़ता है। \emph{(चक्रिकाओं $\overline D(p, r)$, $p \in S$ पर
        हेली लगाइए।)}
\end{enumerate}

\medskip
\noindent\textbf{भाग IV --- केंद्रबिंदु प्रमेय।} समतल के $n$ बिंदुओं के
किसी परिमित समुच्चय $S$ का \emph{केंद्रबिंदु} ऐसा बिंदु $c$ है (जो $S$ में
हो, यह आवश्यक नहीं) कि $c$ को समाहित करने वाला हर संवृत अर्धसमतल $S$ के कम
से कम $n/3$ बिंदु समाहित करे।
\begin{enumerate}[resume]
  \item (विमा $1$) वास्तविक $x_1 \leq \dots \leq x_n$ के लिए दिखाइए कि माध्यिका $c = x_{\lceil n/2 \rceil}$ यह पूरा
        करती है: $c$ को समाहित करने वाली हर संवृत अर्धरेखा $x_i$ में से कम
        से कम $n/2$ समाहित करती है।
  \item (गिनती प्रमेयिका) यदि $S$ के उपसमुच्चय $A, B, C$ ऐसे हों कि $\abs A, \abs B, \abs C > \tfrac{2n}3$, तो
        दिखाइए $A \cap
        B \cap C \neq \emptyset$।
  \item मान लीजिए $m = \floor{2n/3} + 1$, और $\mathcal F$ $\operatorname{conv}(T)$, $T \subseteq S$, $\abs T =
        m$ वाले \omterm{def:b2:affine:convex}{उत्तल} कवचों का
        (परिमित) कुल है। दिखाइए कि $\mathcal F$ के किन्हीं तीन सदस्यों का कोई साझा
        बिंदु है, और हेली से उन सबमें साझा कोई बिंदु $c$ निष्कर्ष रूप में
        निकालिए।
  \item सिद्ध कीजिए कि यही $c$ $S$ का केंद्रबिंदु है: अर्थात्
        \emph{केंद्रबिंदु प्रमेय}। \emph{(यदि $c$ से होकर जाने वाले किसी
        संवृत अर्धसमतल में $n/3$ से कम बिंदु होते, तो उसका \omterm{def:b2:metric:topology}{विवृत} पूरक $m$
        बिंदुओं का कोई समुच्चय $T$ समाहित करता, और $\operatorname{conv}(T)$ $c$ से बच जाता।)}
  \item तीक्ष्णता: $n = 3k$ लीजिए और किसी बड़े त्रिभुज के शीर्षों पर केंद्रित
        छोटी त्रिज्या $\varepsilon$ की तीन चक्रिकाओं में से हर एक में $k$ बिंदु
        रखिए। दिखाइए कि समतल के \emph{प्रत्येक} बिंदु $c$ के लिए $c$ को
        समाहित करने वाले किसी संवृत अर्धसमतल में $S$ के अधिकतम $n/3$ बिंदु
        होते हैं, अतः अचर $1/3$ को सुधारा नहीं जा सकता। \emph{($c$ से
        चक्रिका-केंद्रों की तीन दिशाओं में से दो का कोण अधिकतम $2\pi/3$ होता
        है।)}
\end{enumerate}

\medskip
\noindent\textbf{भाग V --- युंग की प्रमेय और संश्लेषण।}
\begin{enumerate}[resume]
  \item (त्रिभुज प्रमेयिका) मान लीजिए $P, Q, R$ तीन ऐसे बिंदु हैं जिनकी
        जोड़े-जोड़े में दूरियाँ $\leq 1$ हैं। दिखाइए कि वे त्रिज्या $1/\sqrt3$ की किसी
        संवृत चक्रिका में पड़ते हैं। \emph{(यदि कोई कोण $\geq \pi/2$ हो, तो सबसे
        लंबी भुजा को व्यास मानकर चक्रिका लीजिए और माध्यिका सूत्र $\norm{RM}^2 = \tfrac12\norm{RP}^2 +
        \tfrac12\norm{RQ}^2 - \tfrac14\norm{PQ}^2$ काम
        में लीजिए; और यदि त्रिभुज न्यूनकोण हो, तो उसके सबसे बड़े कोण, जो
        $\intco{\pi/3}{\pi/2}$ में पड़ता है, से परित्रिज्या $a/(2\sin
        \widehat A)$ को परिबद्ध कीजिए।)}
  \item (युंग) निष्कर्ष निकालिए: व्यास $\leq 1$ वाला समतल का हर \omterm{def:b2:metric:compact}{संहत} उपसमुच्चय
        त्रिज्या $1/\sqrt3$ की किसी संवृत चक्रिका में समाहित है।
  \item तीक्ष्णता: भुजा $1$ और गुरुत्व केंद्र $G$ वाले समबाहु त्रिभुज
        $A_1A_2A_3$ के लिए प्रत्येक बिंदु $O$ हेतु लाइब्नित्स सर्वसमिका $\sum_i \norm{\vect{OA_i}}^2 =
        3\norm{\vect{OG}}^2 + \sum_i \norm{\vect{GA_i}}^2$
        सिद्ध कीजिए, और निष्कर्ष निकालिए कि तीनों शीर्षों को समाहित करने
        वाली किसी भी चक्रिका की त्रिज्या $\geq
        1/\sqrt3$ है, और समता केवल परिवृत्त के
        लिए।
  \item ($\R^n$ में हेली) $\R^n$ में हेली की प्रमेय कहिए और सिद्ध कीजिए: यदि
        परिमित \omterm{def:b2:affine:convex}{उत्तल समुच्चय} ऐसे हों कि उनमें से कोई भी $n + 1$ प्रतिच्छेद
        करते हों, तो वे सब प्रतिच्छेद करते हैं। \emph{(रादों की प्रमेयिका
        \cref{exo:b2:affine:7} $n + 2$ समुच्चयों को निपटा देती है; फिर प्रश्न 12 की तरह आगमन
        कीजिए।)}
  \item संश्लेषण। शृंखला
        \[
        \text{आनत परतंत्रता} \Rightarrow \text{रादों}
        \Rightarrow \text{हेली} \Rightarrow
        \text{केंद्रबिंदु और युंग},
        \]
        जोड़िए और एक-एक वाक्य में बताइए: रैखिक बीजगणित कहाँ प्रवेश करता
        है, भारों के चिह्न कहाँ, उत्तलता कहाँ, और किस अकेले पग ने समतल की
        विमा का उपयोग किया। अचर $3$ (हेली में), $1/3$ (केंद्रबिंदु) और
        $1/\sqrt3$ (युंग) $\R^n$ में क्या बन जाते हैं? (बिना उपपत्ति के कहिए।)
\end{enumerate}
\end{problem}
